\clearpage
\subsection{一个谱序列}\label{sec:18l}
设 $\symbfit{\gamma}$ 为消失路径的特选基，$\symbfit{\delta}$ 为对它施加与\eref{5.25}中的元素 $\Delta^{1/2}\in\operatorname{Br}_m$ 相应的 Hurwitz 移动后得到的基。注意，虽然这是辫群的正元素，但按\secref{16d}的讨论，在几何上，相应圆盘微分同胚是沿一条平行于边界的圆周所作的负 Dehn 扭转的平方根。\znote{289}{这里容易读反：$\Delta^{1/2}\in\operatorname{Br}_m$ 在辫群中是正元素，而它对应的圆盘微分同胚却是「负」Dehn 扭转的平方根。原因在于辫群与映射类的关系是按几何作用（对消失路径的 Hurwitz 移动）定义的（§16d），其符号约定与「正扭转」的字面直觉相反：整圈 $\Delta$ 对应沿边界平行圆的负 Dehn 扭转，故其一半才记为 $\Delta^{1/2}$。这一支在本书中的作用是：$\symbfit{\delta}$ 由 $\symbfit{\gamma}$ 经这一整族移动得到，于是可借\eref{18.30}逐次算出 $\G$ 在新基上的取值，最后再把基点移到 $-*$ 得到 $\symbfit{\gamma}^!$。[来源：原书 §16d、§18l（印刷页 289、290）]}此外，还要考虑将基点从 $*$ 沿顺时针半圆移到 $-*$，由 $\symbfit{\delta}$ 得到的路径集 $\symbfit{\gamma}^!$；见\sr{fig:18.4}{图~18.4}。

\SourcePage{298}
\SourceFigure{figure-18-4}{0.328153493}{18.4}
设 $L_j^\#$，$j=0,1$，为包含在 $E^{\mathrm{shrink}}$ 中的两个膜，$\widetilde L_j^\#$ 为相应的 (U) 型等变膜。将它们视为 $D\F(\pi)$ 的对象，对其应用\stref{命题}{5.17}。所得谱序列收敛到
\[
HF^*(\widetilde L_0^\#,\widetilde L_1^\#)^{\Z/2}\cong HF^*(L_0^\#,L_1^\#),
\]
其中的同构是\stref{命题}{18.13}的一部分。谱序列的 $E_1$ 页由两类 Floer 群的张量积构成。其中一类是
\begin{equation}\label{eq:18.31}
HF^*(\widetilde\Delta_k^\#,\widetilde L_1^\#)^{\Z/2}\cong HF^*(\Delta_k^\#,L_1^\#),
\end{equation}
其中 $\Delta_k$ 为 $\gamma_k$ 对应的 Lefschetz 顶针，$\widetilde\Delta_k$ 为其在 $\widetilde E$ 中的原像。照常将后者赋予 (B) 型等变膜结构，分次任意选取，使\eref{18.31}左端成为 $H\F(\pi)$ 中的态射空间；再应用\stref{笔记}{18.11}，将其等同于 $E$ 本身中适当的 Floer 上同调群，即右端。出现的另一类 Floer 群可类似地写成
\[
HF^*(\widetilde L_0^\#,\widetilde\Pi_k^\#)^{\Z/2}
\cong HF^*(L_0^\#,\Pi_k^\#)
\cong HF^*(L_0^\#,\Delta_k^{!,\#}),
\]
其中 $\Pi_k,\Delta_k^!$ 为 $\delta_k,\gamma_k^!$ 对应的 Lefschetz 顶针。第二个同构由这些顶针之间显然的同痕给出；该同痕移动边界，但仍保持 Floer 上同调群，因为 Floer 方程的相关解从不到达边界。准确地说，应为 $\Delta_k^!$ 使用与变异描述相容的分次：它可由 $\Delta_{m+1-k}\cap\Delta_k^!$ 的唯一交点具有零 Maslov 指标这一性质刻画，因为这是\secref{5k}中 Koszul 对偶集的一般性质之一。据此，有：

\SourcePage{299}
\StatementNumber{27}\begin{corollary}\label{st:18.27}
存在收敛到 $HF^*(L_0^\#,L_1^\#)$ 的谱序列，其起始页为
\begin{equation}\label{eq:18.32}
E_1^{rs}=\bigl(HF^*(L_0^\#,\Delta_{r+1}^{!,\#})\otimes HF^*(\Delta_{m-r}^\#,L_1^\#)\bigr)^{r+s}.
\end{equation}
\end{corollary}

即使 $L_0^\#=L_1^\#=L^\#$，这一结果也很有意义，此时谱序列的目标化为 $H^*(L;\K)$。注意，\stref{命题}{5.17}也描述了\eref{18.32}上第一个微分的形式：它由两个对象集中相继对象之间互相对偶的态射群中的元素作乘法给出。此外，由\stref{命题}{18.14}，这些群可用消失圈写成 $HF^*(V_k^\#,V_{k+1}^\#)$ 与 $HF^*(V_{m-k}^{!,\#},V_{m+1-k}^{!,\#})$。\znote{291}{关于\eref{18.32}的记号：$r,s$ 是双分次，外层上标 $r+s$ 不是矩阵元的下标，而是取该张量积的总次数 $r+s$ 分量（$HF^*$ 自带 $\Z$-分次），把齐次分量写成上标是全书惯用写法。两个因子的指标之和恒为 $m+1$：$r+1$ 与 $m-r$ 恰好配对，核对下标时很有用。又注意 $\Delta^{!,\#}$ 上标的 $!$ 标记与变异描述配套的那族路径 $\gamma_k^!$（见 §5k 的 Koszul 对偶集），不是求导，也不是取对偶空间。[来源：原书 §18l（印刷页 291）]}

\StatementNumber{28}\begin{note}\label{st:18.28}
最初是 Donaldson 向作者提出了谱序列\eref{18.32}应当存在。他的观点基于对角线的几何分解：为简便，假设 $\pi$ 的任意两个临界值的虚部不同。取对角线 $\Delta\subset E\times E$，其中第一个因子的辛形式变号；在第一个因子上施加 $\im(\pi)$ 的 Hamilton 流，使该对角线变形。对角线的一部分会流出 $E$，但其余部分会逼近稳定流形与不稳定流形的乘积，比较如\bcite{73}。这些流形是沿水平方向走向 $\pm\infty$ 的路径对应的 Lefschetz 顶针，应将其看成\sr{fig:18.4}{图~18.4} 中路径 $\gamma_k,\gamma_k^!$ 的变形版本。在我们的证明中，这一思想被 Beilinson 意义下对角线的代数分解所取代。
\end{note}
